%META {"question":"Quan sát hình: đường thẳng a chứa ba điểm M, A, N (A ở giữa). Nhận xét nào ĐÚNG?","options":["A nằm giữa M và N nên MA + AN = MN","Tia NM và tia NA là hai tia đối nhau","Đoạn thẳng MA và đoạn thẳng AN không có gì chung","Điểm N nằm giữa M và A"],"answer":0,"note":"A nằm giữa hai điểm M, N trên cùng đường thẳng nên MA + AN = MN; tia NM và tia NA có gốc khác nhau nên không là tia đối nhau (tia đối phải chung gốc)."}
\documentclass[border=5pt]{standalone}
\usepackage{tkz-euclide}
% Màu dùng chung cho toàn bộ hình vẽ (ảnh stack với nền tối của web)
\definecolor{tminc}{RGB}{226,232,240}   % chữ + nét chính
\definecolor{tmblue}{RGB}{0,210,255}    % nhấn neon xanh
\definecolor{tmpink}{RGB}{255,0,200}    % nhấn neon hồng
\definecolor{tmgreen}{RGB}{0,255,136}   % nhấn neon xanh lá
\color{tminc}

\begin{document}
\begin{tikzpicture}[line width=1.1pt]
  \coordinate (M) at (0,0);
  \coordinate (A) at (1.7,0);
  \coordinate (N) at (3.4,0);
  \tkzDrawLine[tmblue,add=0.7 and 0.7](M,N)
  \tkzDrawPoints[size=4.5,color=tmpoint,fill=tmpoint](M,A,N)
  \tkzLabelPoint[below left=1pt](M){$M$}
  \tkzLabelPoint[below=3pt](A){$A$}
  \tkzLabelPoint[below right=1pt](N){$N$}
  \node[tminc] at (4.05,0.28) {$a$};
\end{tikzpicture}
\end{document}
